%%%% CAPÍTULO 1 - INTRODUÇÃO
%%
%% Conteúdo extraído VERBATIM de TCC_PARKON_V3.pdf (página 9).
%% Fragmento LaTeX: incluído via \include em tese.tex (sem preâmbulo).
%% Codificação UTF-8.
%%
%% Citações numéricas do original mapeadas para chaves biblatex
%% (padrão {sobrenomePrimeiroAutor}{ano}); serão reconciliadas na tarefa 7.2:
%%   [3] -> \cite{miranda2011}  (MIRANDA, H. F. et al. Mobilidade urbana sustentável e o caso de Curitiba)
%%   [9] -> \cite{wang2011}     (WANG, H. A Reservation-based Smart Parking System)
%% Referências cruzadas a capítulos convertidas para \autoref{} (labels definidos nas tarefas 4.2-4.5).
%% Observação: nenhuma sigla/acrônimo (UFF, API, APK, CLI, RF, RNF) ocorre no texto deste capítulo.

\chapter{Introdução}
\label{cap:introducao}

A crescente urbanização e o aumento da frota de veículos têm intensificado os desafios relacionados à mobilidade urbana~\cite{miranda2011}, especialmente no que se refere à disponibilidade de vagas de estacionamento. Em grandes centros, motoristas enfrentam dificuldades para localizar espaços livres, o que resulta em perda de tempo, aumento do tráfego e maior emissão de poluentes. Nesse cenário, soluções tecnológicas voltadas para otimizar esse processo tornam-se essenciais para melhorar a experiência do usuário e reduzir impactos negativos.

Este trabalho apresenta o desenvolvimento e a análise da plataforma Parkon, cujo objetivo é facilitar a busca, reserva e avaliação de vagas de estacionamento em áreas urbanas. O Parkon é composto por um aplicativo mobile voltado aos motoristas e um portal web administrativo para gestores de estacionamento. O aplicativo permite que os usuários localizem vagas disponíveis em tempo real, realizem reservas antecipadas e recebam rotas para o estacionamento selecionado. Além disso, oferece um sistema de avaliação da qualidade do serviço e da vaga utilizada, promovendo uma experiência mais satisfatória e colaborativa. O portal web permite que administradores cadastrem e gerenciem seus estacionamentos, configurem vagas e acompanhem reservas.

A proposta do Parkon surge da necessidade de reduzir o tempo gasto na procura por vagas e aumentar a conveniência para os motoristas, alinhando-se à tendência de integração entre mobilidade e tecnologia. O desenvolvimento da plataforma foi orientado pela crescente demanda por soluções digitais que otimizem processos cotidianos, especialmente em grandes centros urbanos, onde a falta de vagas é um problema recorrente~\cite{wang2011}.

Este estudo aborda as principais funcionalidades da plataforma, os desafios enfrentados durante sua concepção e as perspectivas para implementação prática. Além disso, discute os benefícios esperados para os usuários e para a mobilidade urbana, considerando aspectos como eficiência, praticidade e sustentabilidade.

O trabalho está estruturado da seguinte forma: o~\autoref{cap:referencial} apresenta o referencial teórico sobre mobilidade urbana e soluções tecnológicas para estacionamento; o~\autoref{cap:proposta} detalha a especificação do sistema, incluindo requisitos e casos de uso; o~\autoref{cap:resultados} apresenta o projeto, protótipo e portal web do sistema; por fim, o~\autoref{cap:conclusao} apresenta as conclusões e sugestões para trabalhos futuros.
